\section{结论与展望}

\subsection{主要发现}

1. \textbf{任务一（采用意愿分类）}：CatBoost调参方案验证集ACC=0.8633，5折CV ACC=0.8730$\pm$0.0025。认知与态度维度构成采用意愿的核心驱动力——knowledge\_anxiety\_gap与green\_tech\_score两个业务交互特征合计贡献41.7\%的重要性，超过所有单一原始特征。类别权重消融结果确认，在准确率唯一评估指标下，不加权方案表现最优。Medium类（误判率26.6\%）是三类中区分难度最高的过渡类别，推测其原因在于该类别在特征空间上位于High与Low之间的模糊地带。

2. \textbf{任务二（家充潜力预测）}：严格排除下游变量后，所有模型均退化为多数类基线（ACC=0.6498）。EDV诊断指向同一根源：住房类型、物业政策、停车条件等对家庭充电条件具决定性影响的变量均未纳入数据集，可观测特征与目标变量的单变量相关系数全部在0.013以下。S1诊断上界（加入真实采用意愿后ACC=0.6694，仅优于基线2.0个百分点）进一步印证，在现有特征集合的约束下，该问题的预测上限极为有限。对负结果的完整记录为相关研究提供了一个可参照的方法论诊断流程。

3. \textbf{任务三（里程焦虑识别）}：Main方案（不含ev\_adoption\_likelihood）验证集ACC=0.8550，AUC=0.9042。全部CV参数组性能高度集中（极差0.00059），信号稳定性明确。边界样本分析是本任务中最具诊断价值的环节——极端评分样本近乎完美的分类准确率表明模型捕获了真实信号；边界样本（评分5/6）的较高误判率反映的是标签二值化在阈值邻域内的固有不确定性，而非模型能力的上限。特征重要性分析中，environmental\_awareness\_score与technology\_affinity\_score合计占据85\%重要性，认知维度对里程焦虑的预测贡献远超其他变量。

\subsection{方法学贡献}

以下几项设计在方法层面具备跨问题的参考意义：

\begin{enumerate}[itemsep=0pt]
\item \textbf{统一CatBoost建模框架}：三项任务在模型选型、评估流程（80/20分层划分 + 5折CV + 独立验证集一次性评估）、数据预处理（共享processed\_data.xlsx）上保持完全一致，为横向比较提供了干净的实验基线和连贯的论文逻辑。
\item \textbf{EDV先行诊断方法}：在投入大规模调参之前，通过多模型对照快速检验特征来源是否携带有效信号。任务二的实践表明，这一流程可以有效避免对低信号数据的无效计算。
\item \textbf{跨任务无泄漏堆叠设计}：S2方案的Task1-no-home模型配合5折OOF概率生成流程，为跨任务特征工程提供了一种严格无泄漏的实现范式。
\item \textbf{消融实验与因果推理的组合使用}：任务三的Main与Ablation对比揭示了一项重要区别——Ablation的统计指标在本地表现更好，但其因果方向决定它不适合生产部署。这一设计思路对存在因果依赖链的多变量预测问题具有潜在的参考价值。
\end{enumerate}

\subsection{局限性}

1. \textbf{超参数搜索规模有限}：三项任务的随机搜索均仅覆盖8组参数组合。后续研究可考虑采用Optuna贝叶斯优化或NNI等自动化调参框架，在更大范围内逼近更优参数配置。

2. \textbf{任务一LR反超CatBoost}：逻辑回归默认权重（ACC=0.8668）略优于调参后的CatBoost（ACC=0.8633）。这一结果有两种可能解释：CatBoost超参数空间尚未被充分遍历，或在当前40维特征空间中，非线性增益已接近上限。

3. \textbf{深度学习模型未覆盖}：基于Transformer的表格数据模型（TabTransformer、FT-Transformer）未在本实验中测试。自注意力机制在特定条件下可能捕捉CatBoost遗漏的高阶特征交互，该可能性有待进一步验证。

4. \textbf{任务二关键特征缺失}：数据集中不包含住房类型、物业政策等对家庭充电条件具决定性影响的变量。若后续能够补充这些字段，任务二的建模结果存在重新评估的空间。

5. \textbf{阈值搜索的边界条件}：任务三类别分布较为均衡，OOF阈值优化在本实验中未见实质提升（最优阈值0.45相较默认0.50，ACC差异不足0.0005）。该步骤作为分析流程的一部分予以保留，以维护论文方法论框架的完整性。
